\documentclass[10pt]{article}
\usepackage[paperwidth=384pt,paperheight=900pt,margin=10pt]{geometry}
\usepackage[T1]{fontenc}
\usepackage{lmodern}
\usepackage[utf8]{inputenc}
\usepackage{amsmath,amssymb}
\usepackage{array}
\usepackage{enumitem}
\usepackage{xcolor}
\usepackage{microtype}

\setlength{\parindent}{0pt}
\setlength{\parskip}{4pt}
\setlist[itemize]{leftmargin=12pt}
\setlist[enumerate]{leftmargin=14pt}

\sloppy
\allowdisplaybreaks
\emergencystretch=2em

\begin{document}

\section*{Uppgift 1}

\subsection*{Original question}
En kruka innehåller tio röda och tio vita bollar.
Man väljer 5 bollar slumpmässigt, utan återläggning.
Låt $A$ vara händelsen att man bara får röda bollar och låt
$B$ vara händelsen att man får minst en boll av vardera färg.
Beräkna följande sannolikheter:

\begin{enumerate}[label=\alph*)]
\item $P(A)$. \hfill (1p)
\item $P(B)$. \hfill (2p)
\item $P(A\cup B)$. \hfill (1p)
\item $P(A\mid B^c)$. \hfill (2p)
\end{enumerate}

\textit{Kommentar:}
Om du ej har löst a), b) så kan du använda värdena
$p(A)=0.1$, $p(B)=0.8$ i c) och d).
Notera dock att detta inte är de korrekta värdena.

\subsection*{Compiled notes from Yacoub + Obid}

\subsubsection*{What we are doing}
We draw 5 balls without replacement from 20 balls total:
10 red and 10 white.

The notes use combinations because order does not matter.

\[
\binom{n}{k}
\]

means the number of ways to choose $k$ objects from $n$ objects.

\subsubsection*{Given information}
\begin{itemize}
\item There are 10 red balls.
\item There are 10 white balls.
\item Total balls:
\[
20
\]
\item We choose:
\[
n=5
\]
\item No return, so this is without replacement.
\end{itemize}

The events are:

\[
A=\text{only red}
\]

This means:

\[
A=\text{5 red and 0 white}.
\]

Also:

\[
B=\text{at least one of each color}.
\]

So:

\[
B=\text{at least 1 red and at least 1 white}.
\]

\subsubsection*{Step-by-step solution}

\textbf{Total number of choices}

The notes first find the total number of ways to choose 5 balls
from 20:

\begin{align*}
\binom{20}{5}
&=
\frac{20\cdot 19\cdot 18\cdot 17\cdot 16}
{5\cdot 4\cdot 3\cdot 2\cdot 1} \\[2mm]
&=15504
\end{align*}

\textbf{a) Calculate $P(A)$}

Event $A$ means all 5 chosen balls are red.

Number of ways to choose 5 red from 10 red:

\begin{align*}
\binom{10}{5}
&=
\frac{10\cdot 9\cdot 8\cdot 7\cdot 6}
{5\cdot 4\cdot 3\cdot 2\cdot 1} \\[2mm]
&=252
\end{align*}

Number of ways to choose 0 white from 10 white:

\[
\binom{10}{0}=1.
\]

So:

\begin{align*}
P(A)
&=
\frac{
\binom{10}{5}\binom{10}{0}
}
{\binom{20}{5}} \\[2mm]
&=
\frac{252}{15504} \\[2mm]
&=
\frac{21}{1292} \\[2mm]
&\approx 0.0163
\end{align*}

So the chance to get 5 red is about:

\[
1.63\,\%.
\]

\textbf{b) Calculate $P(B)$}

En kruka innehåller tio röda och tio vita bollar.
Man väljer 5 bollar slumpmässigt, utan återläggning.
Låt $A$ vara händelsen att man bara får röda bollar och låt
$B$ vara händelsen att man får minst en boll av vardera färg.
Beräkna följande sannolikheter:

\begin{enumerate}[label=\alph*)]
\item $P(A)$. \hfill (1p)
\item $P(B)$. \hfill (2p)
\item $P(A\cup B)$. \hfill (1p)
\item $P(A\mid B^c)$. \hfill (2p)
\end{enumerate}

\textit{Kommentar:}
Om du ej har löst a), b) så kan du använda värdena
$p(A)=0.1$, $p(B)=0.8$ i c) och d).
Notera dock att detta inte är de korrekta värdena.

Event $B$ means at least one ball of each color.

The notes use the complement:

\[
B^c=\text{not at least one of each color}.
\]

That means all chosen balls are one color only:

\[
B^c=(\text{all red})\cup(\text{all white}).
\]

All red:

\[
252
\]

All white is symmetric:

\[
252
\]

So:

\begin{align*}
P(B^c)
&=
\frac{252+252}{15504} \\[2mm]
&=
\frac{504}{15504} \\[2mm]
&=
\frac{21}{646} \\[2mm]
&\approx 0.0325
\end{align*}

Therefore:

\begin{align*}
P(B)
&=
1-P(B^c) \\[2mm]
&=
1-\frac{504}{15504} \\[2mm]
&=
\frac{15000}{15504} \\[2mm]
&=
\frac{625}{646} \\[2mm]
&\approx 0.9675
\end{align*}

\textbf{c) Calculate $P(A\cup B)$}

En kruka innehåller tio röda och tio vita bollar.
Man väljer 5 bollar slumpmässigt, utan återläggning.
Låt $A$ vara händelsen att man bara får röda bollar och låt
$B$ vara händelsen att man får minst en boll av vardera färg.
Beräkna följande sannolikheter:

\begin{enumerate}[label=\alph*)]
\item $P(A)$. \hfill (1p)
\item $P(B)$. \hfill (2p)
\item $P(A\cup B)$. \hfill (1p)
\item $P(A\mid B^c)$. \hfill (2p)
\end{enumerate}

\textit{Kommentar:}
Om du ej har löst a), b) så kan du använda värdena
$p(A)=0.1$, $p(B)=0.8$ i c) och d).
Notera dock att detta inte är de korrekta värdena.

The notes use:

\[
P(A\cup B)
=P(A)+P(B)-P(A\cap B).
\]

But $A$ and $B$ cannot happen at the same time:

\begin{itemize}
\item $A$ means all red, so no white is allowed.
\item $B$ requires at least one white.
\end{itemize}

So:

\[
P(A\cap B)=0.
\]

Therefore:

\begin{align*}
P(A\cup B)
&=
P(A)+P(B) \\[2mm]
&=
\frac{252}{15504}
+\frac{15000}{15504} \\[2mm]
&=
\frac{15252}{15504} \\[2mm]
&=
\frac{1271}{1292} \\[2mm]
&\approx 0.9837
\end{align*}

\textbf{d) Calculate $P(A\mid B^c)$}

En kruka innehåller tio röda och tio vita bollar.
Man väljer 5 bollar slumpmässigt, utan återläggning.
Låt $A$ vara händelsen att man bara får röda bollar och låt
$B$ vara händelsen att man får minst en boll av vardera färg.
Beräkna följande sannolikheter:

\begin{enumerate}[label=\alph*)]
\item $P(A)$. \hfill (1p)
\item $P(B)$. \hfill (2p)
\item $P(A\cup B)$. \hfill (1p)
\item $P(A\mid B^c)$. \hfill (2p)
\end{enumerate}

\textit{Kommentar:}
Om du ej har löst a), b) så kan du använda värdena
$p(A)=0.1$, $p(B)=0.8$ i c) och d).
Notera dock att detta inte är de korrekta värdena.

The conditional probability formula is:

\begin{align*}
P(A\mid B^c)
&=
\frac{P(A\cap B^c)}
{P(B^c)}.
\end{align*}

If $A$ happens, then all balls are red.
That automatically means $B$ failed, since $B$ needs at least
one white.

So:

\[
A\cap B^c=A.
\]

Therefore:

\begin{align*}
P(A\mid B^c)
&=
\frac{P(A)}{P(B^c)} \\[2mm]
&=
\frac{252/15504}{504/15504} \\[2mm]
&=
\frac{252}{504} \\[2mm]
&=
\frac{1}{2} \\[2mm]
&=0.5
\end{align*}

\subsubsection*{Final answer}
\begin{enumerate}[label=\alph*)]
\item
\[
P(A)=\frac{252}{15504}
=\frac{21}{1292}
\approx 0.0163.
\]

\item
\[
P(B)=\frac{15000}{15504}
=\frac{625}{646}
\approx 0.9675.
\]

\item
\[
P(A\cup B)
=\frac{15252}{15504}
=\frac{1271}{1292}
\approx 0.9837.
\]

\item
\[
P(A\mid B^c)=\frac{1}{2}=0.5.
\]
\end{enumerate}

\subsubsection*{What to remember}
When choosing without replacement and order does not matter,
use combinations:

\[
\binom{n}{k}.
\]

For ``at least one of each color'', it can be easier to use
the complement:

\[
B^c=\text{all red or all white}.
\]

For conditional probability:

\[
P(A\mid B^c)
=
\frac{P(A\cap B^c)}{P(B^c)}.
\]

\section*{Uppgift 2}

\subsection*{Original question}
Betrakta funktionen

\[
f(x)=
\begin{cases}
\dfrac{a}{\sqrt{x}},
& \text{om } 0<x\leq 4,\\
0,
& \text{annars.}
\end{cases}
\]

\begin{enumerate}[label=\alph*)]
\item
Bestäm $a$ så att $f$ bildar en täthetsfunktion för en stokastisk
variabel $X$.
\hfill (1.5p)

\item
Bestäm fördelningsfunktionen $F_X$.
\hfill (1.5p)

\item
Beräkna $P(X\leq 1)$.
\hfill (1p)

\item
Beräkna väntevärde och standardavvikelse för $X$.
\hfill (2p)
\end{enumerate}

\subsection*{Compiled notes from Yacoub + Obid}

\subsubsection*{What we are doing}
The notes treat $f$ as a density function.
First we choose $a$ so that the total area under the PDF is $1$.

Then we integrate the PDF from $0$ to $x$ to get the CDF.
After that we calculate $P(X\leq 1)$ using the CDF.

Finally, we calculate expected value and standard deviation.

\subsubsection*{Given information}
The function is:

\[
f(x)=
\frac{a}{\sqrt{x}},
\qquad 0<x\leq 4.
\]

The notes rewrite:

\[
\frac{1}{\sqrt{x}}=x^{-1/2}.
\]

\subsubsection*{Step-by-step solution}

\textbf{a) Find $a$}

For a PDF:

\[
\int_0^4 \frac{a}{\sqrt{x}}\,dx=1.
\]

Rewrite:

\begin{align*}
a\int_0^4 x^{-1/2}\,dx
&=1
\end{align*}

Integrate:

\begin{align*}
a\left[
2\sqrt{x}
\right]_0^4
&=1
\end{align*}

Plug in the limits:

\begin{align*}
a(2\sqrt{4}-2\sqrt{0})
&=1 \\[2mm]
a(4)
&=1 \\[2mm]
4a
&=1 \\[2mm]
a
&=\frac{1}{4}
\end{align*}

So the full PDF is:

\[
f(x)=\frac{1}{4\sqrt{x}},
\qquad 0<x\leq 4.
\]

\textbf{b) Find the CDF}

Betrakta funktionen

\[
f(x)=
\begin{cases}
\dfrac{a}{\sqrt{x}},
& \text{om } 0<x\leq 4,\\
0,
& \text{annars.}
\end{cases}
\]

\begin{enumerate}[label=\alph*)]
\item
Bestäm $a$ så att $f$ bildar en täthetsfunktion för en stokastisk
variabel $X$.
\hfill (1.5p)

\item
Bestäm fördelningsfunktionen $F_X$.
\hfill (1.5p)

\item
Beräkna $P(X\leq 1)$.
\hfill (1p)

\item
Beräkna väntevärde och standardavvikelse för $X$.
\hfill (2p)
\end{enumerate}

The CDF is:

\[
F_X(x)=P(X\leq x).
\]

For $0<x\leq 4$:

\begin{align*}
F_X(x)
&=
\int_0^x
\frac{1}{4\sqrt{t}}\,dt \\[2mm]
&=
\frac{1}{4}
\int_0^x
t^{-1/2}\,dt \\[2mm]
&=
\frac{1}{4}
\left[
2\sqrt{t}
\right]_0^x \\[2mm]
&=
\frac{1}{4}
(2\sqrt{x}-0) \\[2mm]
&=
\frac{1}{2}\sqrt{x}
\end{align*}

So:

\[
F_X(x)=
\begin{cases}
0,
& x\leq 0,\\
\dfrac{1}{2}\sqrt{x},
& 0<x\leq 4,\\
1,
& x>4.
\end{cases}
\]

The notes also remind us:
the CDF cannot go bigger than $1$.

\textbf{c) Calculate $P(X\leq 1)$}

Betrakta funktionen

\[
f(x)=
\begin{cases}
\dfrac{a}{\sqrt{x}},
& \text{om } 0<x\leq 4,\\
0,
& \text{annars.}
\end{cases}
\]

\begin{enumerate}[label=\alph*)]
\item
Bestäm $a$ så att $f$ bildar en täthetsfunktion för en stokastisk
variabel $X$.
\hfill (1.5p)

\item
Bestäm fördelningsfunktionen $F_X$.
\hfill (1.5p)

\item
Beräkna $P(X\leq 1)$.
\hfill (1p)

\item
Beräkna väntevärde och standardavvikelse för $X$.
\hfill (2p)
\end{enumerate}

Use the CDF:

\begin{align*}
P(X\leq 1)
&=
F_X(1) \\[2mm]
&=
\frac{1}{2}\sqrt{1} \\[2mm]
&=
\frac{1}{2}
\end{align*}

\textbf{d) Expected value and standard deviation}

Betrakta funktionen

\[
f(x)=
\begin{cases}
\dfrac{a}{\sqrt{x}},
& \text{om } 0<x\leq 4,\\
0,
& \text{annars.}
\end{cases}
\]

\begin{enumerate}[label=\alph*)]
\item
Bestäm $a$ så att $f$ bildar en täthetsfunktion för en stokastisk
variabel $X$.
\hfill (1.5p)

\item
Bestäm fördelningsfunktionen $F_X$.
\hfill (1.5p)

\item
Beräkna $P(X\leq 1)$.
\hfill (1p)

\item
Beräkna väntevärde och standardavvikelse för $X$.
\hfill (2p)
\end{enumerate}

\textbf{Find $E[X]$}

The notes use:

\[
E[X]=\int x f(x)\,dx.
\]

So:

\begin{align*}
E[X]
&=
\int_0^4
x\left(\frac{1}{4\sqrt{x}}\right)\,dx \\[2mm]
&=
\int_0^4
\frac{x}{4\sqrt{x}}\,dx
\end{align*}

Since

\[
\frac{x}{\sqrt{x}}=x^{1/2},
\]

we get:

\begin{align*}
E[X]
&=
\frac{1}{4}
\int_0^4 x^{1/2}\,dx \\[2mm]
&=
\frac{1}{4}
\left[
\frac{x^{3/2}}{3/2}
\right]_0^4 \\[2mm]
&=
\frac{1}{4}
\left(
\frac{4^{3/2}}{3/2}-0
\right) \\[2mm]
&=
\frac{1}{4}\cdot \frac{16}{3} \\[2mm]
&=
\frac{4}{3}
\end{align*}

\textbf{Find $E[X^2]$}

The notes use:

\[
E[X^2]
=
\int x^2 f(x)\,dx.
\]

So:

\begin{align*}
E[X^2]
&=
\int_0^4
x^2\left(\frac{1}{4}x^{-1/2}\right)\,dx \\[2mm]
&=
\frac{1}{4}
\int_0^4
x^{3/2}\,dx \\[2mm]
&=
\frac{1}{4}
\left[
\frac{x^{5/2}}{5/2}
\right]_0^4 \\[2mm]
&=
\frac{1}{4}
\left(
\frac{4^{5/2}}{5/2}
-0
\right) \\[2mm]
&=
\frac{1}{4}
\left(
\frac{32}{5/2}
\right) \\[2mm]
&=
\frac{1}{4}\cdot \frac{64}{5} \\[2mm]
&=
\frac{16}{5}
\end{align*}

Variance:

\begin{align*}
\operatorname{Var}(X)
&=
E[X^2]-(E[X])^2 \\[2mm]
&=
\frac{16}{5}
-\left(\frac{4}{3}\right)^2 \\[2mm]
&=
\frac{16}{5}
-\frac{16}{9} \\[2mm]
&=
\frac{64}{45}
\end{align*}

Standard deviation:

\begin{align*}
\sigma
&=
\sqrt{\operatorname{Var}(X)} \\[2mm]
&=
\sqrt{\frac{64}{45}} \\[2mm]
&=
\frac{8}{3\sqrt{5}}
\end{align*}

\subsubsection*{Final answer}
\begin{enumerate}[label=\alph*)]
\item
\[
a=\frac{1}{4}.
\]

\item
\[
F_X(x)=
\begin{cases}
0,
& x\leq 0,\\
\dfrac{1}{2}\sqrt{x},
& 0<x\leq 4,\\
1,
& x>4.
\end{cases}
\]

\item
\[
P(X\leq 1)=\frac{1}{2}.
\]

\item
\[
E[X]=\frac{4}{3},
\qquad
\sigma=\frac{8}{3\sqrt{5}}.
\]
\end{enumerate}

\subsubsection*{What to remember}
For a PDF:

\[
\int f(x)\,dx=1.
\]

For a CDF:

\[
F_X(x)=\int_{-\infty}^{x} f(t)\,dt.
\]

For expected value:

\[
E[X]=\int x f(x)\,dx.
\]

For standard deviation:

\[
\sigma=
\sqrt{E[X^2]-(E[X])^2}.
\]

\section*{Uppgift 3}

\subsection*{Original question}
Låt \(X \in \operatorname{Exp}(2)\),
\(Y \in \operatorname{Exp}(3)\) och antag att
\(X, Y\) är oberoende s.v.

\begin{enumerate}[label=\alph*)]
\item Beräkna \(P(1 \leq X < 3)\). \hfill (1p)
\item Beräkna \(E(3X+Y)\) och \(V(3X+Y)\). \hfill (2p)
\item Bestäm fördelningsfunktionen för \(4X+1\). \hfill (1.5p)
\item Beräkna \(p(\max\{X,Y\}\leq 1)\). \hfill (1.5p)
\end{enumerate}

\subsection*{Compiled notes from Yacoub + Obid}

\subsubsection*{What we are doing}
The notes use the exponential distribution.
For \(X\) and \(Y\), we use the CDF,
the mean, and the variance.

\subsubsection*{Given information}
\[
X\in \operatorname{Exp}(2)
\]
\[
Y\in \operatorname{Exp}(3)
\]

\[
X \text{ and } Y
\text{ are independent.}
\]

For an exponential distribution:

\[
F(x)=1-e^{-\lambda x},
\qquad x\geq 0.
\]

\[
E[X]=\frac{1}{\lambda}
\]

\[
\operatorname{Var}(X)=\frac{1}{\lambda^2}
\]

So the notes write:

\[
X\in \operatorname{Exp}(2)
\Rightarrow
E[X]=\frac{1}{2},
\quad
\operatorname{Var}(X)=\frac{1}{4}.
\]

\[
Y\in \operatorname{Exp}(3)
\Rightarrow
E[Y]=\frac{1}{3},
\quad
\operatorname{Var}(Y)=\frac{1}{9}.
\]

\subsubsection*{Step-by-step solution}

\textbf{a) Calculate \(P(1\leq X<3)\)}

Use the CDF for \(X\):

\[
F_X(x)=1-e^{-2x}.
\]

The notes calculate:

\begin{align*}
P(1\leq X<3)
&=F_X(3)-F_X(1) \\[2mm]
&=(1-e^{-2\cdot 3})
 -(1-e^{-2\cdot 1}) \\[2mm]
&=(1-e^{-6})-(1-e^{-2}) \\[2mm]
&=e^{-2}-e^{-6} \\[2mm]
&\approx 0.1328566311.
\end{align*}

\[
P(1\leq X<3)
\approx 0.1329.
\]

\textbf{b) Calculate \(E(3X+Y)\)}

Låt \(X \in \operatorname{Exp}(2)\),
\(Y \in \operatorname{Exp}(3)\) och antag att
\(X, Y\) är oberoende s.v.

\begin{enumerate}[label=\alph*)]
\item Beräkna \(P(1 \leq X < 3)\). \hfill (1p)
\item Beräkna \(E(3X+Y)\) och \(V(3X+Y)\). \hfill (2p)
\item Bestäm fördelningsfunktionen för \(4X+1\). \hfill (1.5p)
\item Beräkna \(p(\max\{X,Y\}\leq 1)\). \hfill (1.5p)
\end{enumerate}

Use the expected values from the notes:

\[
E[X]=\frac{1}{2},
\qquad
E[Y]=\frac{1}{3}.
\]

\begin{align*}
E(3X+Y)
&=3E[X]+E[Y] \\[2mm]
&=3\cdot \frac{1}{2}
  +\frac{1}{3} \\[2mm]
&=\frac{3}{2}+\frac{1}{3} \\[2mm]
&=\frac{9}{6}+\frac{2}{6} \\[2mm]
&=\frac{11}{6}.
\end{align*}

\textbf{b) Calculate \(V(3X+Y)\)}

Låt \(X \in \operatorname{Exp}(2)\),
\(Y \in \operatorname{Exp}(3)\) och antag att
\(X, Y\) är oberoende s.v.

\begin{enumerate}[label=\alph*)]
\item Beräkna \(P(1 \leq X < 3)\). \hfill (1p)
\item Beräkna \(E(3X+Y)\) och \(V(3X+Y)\). \hfill (2p)
\item Bestäm fördelningsfunktionen för \(4X+1\). \hfill (1.5p)
\item Beräkna \(p(\max\{X,Y\}\leq 1)\). \hfill (1.5p)
\end{enumerate}

Because \(X\) and \(Y\) are independent,
the notes add the variances:

\begin{align*}
\operatorname{Var}(3X+Y)
&=3^2\operatorname{Var}(X)
  +1^2\operatorname{Var}(Y) \\[2mm]
&=9\cdot \frac{1}{4}
  +\frac{1}{9} \\[2mm]
&=\frac{9}{4}+\frac{1}{9} \\[2mm]
&=\frac{81}{36}
  +\frac{4}{36} \\[2mm]
&=\frac{85}{36}.
\end{align*}

\textbf{c) Find the CDF for \(4X+1\)}

Låt \(X \in \operatorname{Exp}(2)\),
\(Y \in \operatorname{Exp}(3)\) och antag att
\(X, Y\) är oberoende s.v.

\begin{enumerate}[label=\alph*)]
\item Beräkna \(P(1 \leq X < 3)\). \hfill (1p)
\item Beräkna \(E(3X+Y)\) och \(V(3X+Y)\). \hfill (2p)
\item Bestäm fördelningsfunktionen för \(4X+1\). \hfill (1.5p)
\item Beräkna \(p(\max\{X,Y\}\leq 1)\). \hfill (1.5p)
\end{enumerate}

The notes set:

\[
Z=4X+1.
\]

We want:

\[
F_Z(z)=P(Z\leq z).
\]

\begin{align*}
F_Z(z)
&=P(4X+1\leq z) \\[2mm]
&=P(4X\leq z-1) \\[2mm]
&=P\left(X\leq \frac{z-1}{4}\right) \\[2mm]
&=F_X\left(\frac{z-1}{4}\right).
\end{align*}

Plug this into the CDF for \(X\):

\begin{align*}
F_Z(z)
&=1-e^{-2\left(\frac{z-1}{4}\right)} \\[2mm]
&=1-e^{-\frac{z-1}{2}}.
\end{align*}

Since \(X\geq 0\), we have:

\[
4X+1\geq 1.
\]

So the CDF is:

\[
F_{4X+1}(z)=
\begin{cases}
0, & z<1,\\[1mm]
1-e^{-\frac{z-1}{2}},
& z\geq 1.
\end{cases}
\]

\textbf{d) Calculate \(p(\max\{X,Y\}\leq 1)\)}

Låt \(X \in \operatorname{Exp}(2)\),
\(Y \in \operatorname{Exp}(3)\) och antag att
\(X, Y\) är oberoende s.v.

\begin{enumerate}[label=\alph*)]
\item Beräkna \(P(1 \leq X < 3)\). \hfill (1p)
\item Beräkna \(E(3X+Y)\) och \(V(3X+Y)\). \hfill (2p)
\item Bestäm fördelningsfunktionen för \(4X+1\). \hfill (1.5p)
\item Beräkna \(p(\max\{X,Y\}\leq 1)\). \hfill (1.5p)
\end{enumerate}

The notes rewrite the event:

\[
\max\{X,Y\}\leq 1
\]

as:

\[
X\leq 1
\quad \text{and} \quad
Y\leq 1.
\]

Because \(X\) and \(Y\) are independent:

\begin{align*}
P(\max\{X,Y\}\leq 1)
&=P(X\leq 1,\;Y\leq 1) \\[2mm]
&=P(X\leq 1)\cdot P(Y\leq 1) \\[2mm]
&=F_X(1)\cdot F_Y(1).
\end{align*}

Plug in the CDF values:

\begin{align*}
F_X(1)
&=1-e^{-2\cdot 1}
=1-e^{-2},
\\[2mm]
F_Y(1)
&=1-e^{-3\cdot 1}
=1-e^{-3}.
\end{align*}

Therefore:

\begin{align*}
P(\max\{X,Y\}\leq 1)
&=(1-e^{-2})(1-e^{-3}).
\end{align*}

\subsubsection*{Final answer}
\begin{enumerate}[label=\alph*)]
\item
\[
P(1\leq X<3)
=e^{-2}-e^{-6}
\approx 0.1329.
\]

\item
\[
E(3X+Y)=\frac{11}{6}.
\]
\[
V(3X+Y)=\frac{85}{36}.
\]

\item
\[
F_{4X+1}(z)=
\begin{cases}
0, & z<1,\\[1mm]
1-e^{-\frac{z-1}{2}},
& z\geq 1.
\end{cases}
\]

\item
\[
P(\max\{X,Y\}\leq 1)
=(1-e^{-2})(1-e^{-3}).
\]
\end{enumerate}

\subsubsection*{What to remember}
For \(X\in \operatorname{Exp}(\lambda)\):

\[
F_X(x)=1-e^{-\lambda x}.
\]

\[
E[X]=\frac{1}{\lambda},
\qquad
\operatorname{Var}(X)=\frac{1}{\lambda^2}.
\]

For independent variables:

\[
\operatorname{Var}(aX+bY)
=a^2\operatorname{Var}(X)
 +b^2\operatorname{Var}(Y).
\]

\[
P(X\leq a,\;Y\leq b)
=P(X\leq a)P(Y\leq b).
\]

\section*{Uppgift 4}

\subsection*{Original question}
% UNCERTAIN: original question image was not provided in this prompt;
% this wording is transcribed from the provided notes.
Låt \(x_1,x_2,\ldots,x_n\) vara ett stickprov
från en fördelning med täthetsfunktion

\[
f_X(x)=
\begin{cases}
\dfrac{x}{a}e^{-\frac{x^2}{2a}},
& x>0,\\[2mm]
0, & \text{annars.}
\end{cases}
\]

Ange ML-skattningen av \(a\).

\subsection*{Compiled notes from Yacoub + Obid}

\subsubsection*{What we are doing}
The notes find the value of \(a\) that makes the
observed data most likely.
That means we build \(L(a)\), take \(\ln\),
differentiate, set equal to zero,
and isolate \(a\).

\subsubsection*{Given information}
The density used in the notes is:

\[
f_X(x_i)
=
\frac{x_i}{a}
e^{-\frac{x_i^2}{2a}}.
\]

We have observations:

\[
x_1,x_2,\ldots,x_n.
\]

\subsubsection*{Step-by-step solution}

\textbf{Step 1: Likelihood function}

The notes start with:

\[
L(a)=
\prod_{i=1}^{n} f_X(x_i).
\]

Substitute the density:

\begin{align*}
L(a)
&=
\prod_{i=1}^{n}
\left(
\frac{x_i}{a}
e^{-\frac{x_i^2}{2a}}
\right).
\end{align*}

\textbf{Step 2: Take logarithm}

The notes use logarithms because they help with
products and exponential functions:

\[
\ln(AB)=\ln A+\ln B.
\]

\[
\ln\left(\prod(\cdot)\right)
=
\sum \ln(\cdot).
\]

So:

\begin{align*}
l(a)
&=\ln(L(a)) \\[2mm]
&=
\sum_{i=1}^{n}
\ln\left(
\frac{x_i}{a}
e^{-\frac{x_i^2}{2a}}
\right).
\end{align*}

Split the logarithm:

\begin{align*}
l(a)
&=
\sum_{i=1}^{n}
\left(
\ln(x_i)-\ln(a)
-\frac{x_i^2}{2a}
\right) \\[2mm]
&=
\sum_{i=1}^{n}\ln(x_i)
-n\ln(a)
-\frac{1}{2a}
\sum_{i=1}^{n}x_i^2.
\end{align*}

\textbf{Step 3: Differentiate with respect to \(a\)}

The notes differentiate \(l(a)\):

\begin{align*}
\frac{d}{da}l(a)
&=
0-\frac{n}{a}
+
\frac{1}{2a^2}
\sum_{i=1}^{n}x_i^2 \\[2mm]
&=
-\frac{n}{a}
+
\frac{1}{2a^2}
\sum_{i=1}^{n}x_i^2.
\end{align*}

\textbf{Step 4: Set derivative equal to zero}

\begin{align*}
-\frac{n}{a}
+
\frac{1}{2a^2}
\sum_{i=1}^{n}x_i^2
&=0.
\end{align*}

Move \(\frac{n}{a}\) to the other side:

\begin{align*}
\frac{1}{2a^2}
\sum_{i=1}^{n}x_i^2
&=
\frac{n}{a}.
\end{align*}

Multiply by \(a^2\):

\begin{align*}
\frac{1}{2}
\sum_{i=1}^{n}x_i^2
&=
na.
\end{align*}

Divide by \(n\):

\begin{align*}
a
&=
\frac{1}{2n}
\sum_{i=1}^{n}x_i^2.
\end{align*}

\subsubsection*{Final answer}
The ML estimate is:

\[
\hat a
=
\frac{1}{2n}
\sum_{i=1}^{n}x_i^2.
\]

\subsubsection*{What to remember}
For ML estimation:

\[
L(\text{parameter})
=
\prod f(x_i).
\]

Then:

\[
l(\text{parameter})
=\ln(L(\text{parameter})).
\]

To find the maximum:

\[
\frac{d}{d(\text{parameter})}l=0.
\]

\section*{Uppgift 5}

\subsection*{Original question}
Låt
\[
13.4,\;14.1,\;15.1,\;14.3,\;14.9,\;13.9,\;14.4,\;14.0
\]
vara ett stickprov från \(N(\mu,\sigma)\).
Beräkning ger att
\[
\sum_{i=1}^{8}x_i=114.1,
\qquad
\sum_{i=1}^{8}x_i^2=1629.45.
\]

\begin{enumerate}[label=\alph*)]
\item Konstruera ett (tvåsidigt) 95\%
konfidensintervall för \(\mu\) under
förutsättning att \(\sigma=0.5\).
\hfill (3p)

\item Konstruera ett (tvåsidigt) 95\%
konfidensintervall för \(\mu\) under
förutsättning att \(\sigma\) är okänt.
\hfill (3p)
\end{enumerate}

\subsection*{Compiled notes from Yacoub + Obid}

\subsubsection*{What we are doing}
The notes construct confidence intervals for the
mean \(\mu\).

In part a), \(\sigma\) is known, so the notes use
a \(z\)-score.

In part b), \(\sigma\) is unknown, so the notes first
calculate \(s\), then use a \(t\)-value.

\subsubsection*{Given information}
\[
n=8
\]

\[
\sum x_i=114.1
\]

\[
\sum x_i^2=1629.45
\]

The sample mean is:

\begin{align*}
\bar{x}
&=\frac{\sum x_i}{n} \\[2mm]
&=\frac{114.1}{8} \\[2mm]
&=14.2625.
\end{align*}

\subsubsection*{Step-by-step solution}

\textbf{a) Known \(\sigma=0.5\)}

For a 95\% confidence interval, the notes split
the remaining 5\% into two tails:

\[
\alpha=0.05,
\qquad
\frac{\alpha}{2}=0.025.
\]

From the normal table:

\[
z_{0.025}=1.96.
\]

The margin of error is:

\[
ME=z\cdot \frac{\sigma}{\sqrt{n}}.
\]

Substitute the values from the notes:

\begin{align*}
ME
&=1.96\cdot \frac{0.5}{\sqrt{8}} \\[2mm]
&\approx 0.3465.
\end{align*}

Construct the interval:

\begin{align*}
\bar{x}\pm ME
&=14.2625\pm 0.3465.
\end{align*}

Lower bound:

\begin{align*}
14.2625-0.3465
&=13.916.
\end{align*}

Upper bound:

\begin{align*}
14.2625+0.3465
&=14.609.
\end{align*}

So:

\[
\mu\in [13.916,\;14.609].
\]

\textbf{Extra expanded notes for a)}

The added notes write the given values again:

\[
n=8
\]

\[
\sum_{i=1}^{8}x_i=114.1
\]

\[
\sum_{i=1}^{8}x_i^2=1629.45
\]

For a 95\% confidence interval:

\[
100\%-95\%=5\%.
\]

Because the interval is two-sided,
split the 5\% into two tails:

\[
\frac{5\%}{2}=2.5\%=0.025.
\]

So the note uses:

\[
z_{0.025}=1.9600.
\]

The margin of error is:

\[
ME=
z\cdot \frac{\sigma}{\sqrt{n}}.
\]

Substitute:

\begin{align*}
ME
&=
1.96\cdot
\frac{0.5}{\sqrt{8}} \\[2mm]
&=0.3465.
\end{align*}

The sample mean is:

\begin{align*}
\bar{x}
&=
\frac{\sum x_i}{n} \\[2mm]
&=
\frac{114.1}{8} \\[2mm]
&=14.2625.
\end{align*}

Then:

\begin{align*}
\bar{x}\pm ME
&=14.2625\pm 0.3465.
\end{align*}

So the expanded note gives:

\[
[13.916,\;14.609].
\]

\textbf{b) Unknown \(\sigma\)}

Låt
\[
13.4,\;14.1,\;15.1,\;14.3,\;14.9,\;13.9,\;14.4,\;14.0
\]
vara ett stickprov från \(N(\mu,\sigma)\).
Beräkning ger att
\[
\sum_{i=1}^{8}x_i=114.1,
\qquad
\sum_{i=1}^{8}x_i^2=1629.45.
\]

\begin{enumerate}[label=\alph*)]
\item Konstruera ett (tvåsidigt) 95\%
konfidensintervall för \(\mu\) under
förutsättning att \(\sigma=0.5\).
\hfill (3p)

\item Konstruera ett (tvåsidigt) 95\%
konfidensintervall för \(\mu\) under
förutsättning att \(\sigma\) är okänt.
\hfill (3p)
\end{enumerate}

When \(\sigma\) is unknown, the notes use the
sample standard deviation \(s\).

Formula from the notes:

\[
s^2=
\frac{1}{n-1}
\left(
\sum x_i^2
-
\frac{(\sum x_i)^2}{n}
\right).
\]

Substitute:

\begin{align*}
s^2
&=
\frac{1}{8-1}
\left(
1629.45
-
\frac{(114.1)^2}{8}
\right) \\[2mm]
&=
\frac{1}{7}
\left(
1629.45
-
\frac{13018.81}{8}
\right) \\[2mm]
&=
\frac{1}{7}
\left(
1629.45-1627.35125
\right) \\[2mm]
&=
\frac{2.09875}{7} \\[2mm]
&\approx 0.29982.
\end{align*}

Then:

\begin{align*}
s
&=\sqrt{s^2} \\[2mm]
&=\sqrt{0.29982} \\[2mm]
&\approx 0.5476.
\end{align*}

Degrees of freedom:

\[
df=n-1=7.
\]

For 95\% confidence:

\[
\alpha=0.05,
\qquad
\frac{\alpha}{2}=0.025.
\]

From the \(t\)-table, the notes use:

\[
t_{0.025,7}=2.365.
\]

Margin of error:

\begin{align*}
ME
&=t\cdot \frac{s}{\sqrt{n}} \\[2mm]
&=2.365\cdot \frac{0.5476}{\sqrt{8}} \\[2mm]
&\approx 0.4579.
\end{align*}

Construct the interval:

\begin{align*}
\bar{x}\pm ME
&=14.2625\pm 0.4579.
\end{align*}

Lower bound:

\begin{align*}
14.2625-0.4579
&=13.8046.
\end{align*}

Upper bound:

\begin{align*}
14.2625+0.4579
&=14.7204.
\end{align*}

So:

\[
\mu\in [13.8046,\;14.7204].
\]

\subsubsection*{Final answer}
\begin{enumerate}[label=\alph*)]
\item If \(\sigma=0.5\), the 95\% confidence interval is:
\[
[13.916,\;14.609].
\]

\item If \(\sigma\) is unknown, the 95\% confidence interval is:
\[
[13.8046,\;14.7204].
\]
\end{enumerate}

\subsubsection*{What to remember}
If \(\sigma\) is known, use:

\[
\bar{x}\pm z\cdot \frac{\sigma}{\sqrt{n}}.
\]

If \(\sigma\) is unknown, calculate \(s\) and use:

\[
\bar{x}\pm t\cdot \frac{s}{\sqrt{n}}.
\]

\[
df=n-1.
\]

\section*{Uppgift 6}

\subsection*{Original question}
En bonde köper 10000 granplantor.
Höjden på dessa plantor antas vara normalfördelade
med väntevärde \(\mu\) cm och standardavvikelse
10 cm.
Säljaren av dessa plantor påstår att den genomsnittliga
höjden på plantorna är \(\mu=70\) cm.
Om \(\mu<70\) så har bonden rätt att underkänna
leveransen av plantor.

\begin{enumerate}[label=\alph*)]
\item Bonden misstänker att \(\mu<70\) cm.
För att testa detta, så mäter bonden 100 plantor
och finner att \(\bar{x}=68\) cm.
Kan bonden underkänna leveransen om han är villig
att ta en risk på 5\% att han gör fel?
Svara på denna fråga genom att ställa upp ett
lämpligt hypotestest (ställ upp hypotes \(H_0\)
och mothypotes \(H_1\)).
\hfill (3p)

\item Antag nu istället att bonden mäter 200 plantor
och finner att \(\bar{x}=69\).
Kan bonden underkänna leveransen om han är villig
att ta en risk på 5\% att han gör fel?
Svara på denna fråga genom att använda samma
hypotestest som du använde i a).
\hfill (3p)
\end{enumerate}

\subsection*{Compiled notes from Yacoub + Obid}

\subsubsection*{What we are doing}
The notes set up a hypothesis test.
The farmer wants to know if the plants are shorter
than \(70\) cm on average.

Since the alternative is \(\mu<70\),
this is a left-tailed test.

\subsubsection*{Given information}
From the question and notes:

\[
\sigma=10
\]

\[
\mu_0=70
\]

\[
\alpha=0.05
\]

For part a):

\[
n=100,
\qquad
\bar{x}=68.
\]

For part b):

\[
n=200,
\qquad
\bar{x}=69.
\]

\subsubsection*{Step-by-step solution}

\textbf{a) Set up the hypotheses}

The notes write:

\[
H_0:\mu=70.
\]

\[
H_1:\mu<70.
\]

Because the sign is \(<\), the test is left-tailed.

\[
\alpha=0.05.
\]

For a left-tailed test at 5\%,
the notes use the critical value:

\[
z_\alpha\approx -1.645.
\]

\textbf{a) Test statistic}

The notes use the \(z\)-test formula:

\[
z=
\frac{\bar{x}-\mu_0}
{\sigma/\sqrt{n}}.
\]

Plug in the values:

\begin{align*}
z
&=
\frac{68-70}
{10/\sqrt{100}} \\[2mm]
&=
\frac{-2}{10/10} \\[2mm]
&=
\frac{-2}{1} \\[2mm]
&=-2.
\end{align*}

Compare with the critical value:

\[
-2<-1.645.
\]

So the test statistic is in the rejection region.

Therefore:

\[
\text{Reject } H_0.
\]

The notes conclude that \(H_1\) is supported:
the farmer can reject the delivery.

\textbf{b) Same hypothesis test}

En bonde köper 10000 granplantor.
Höjden på dessa plantor antas vara normalfördelade
med väntevärde \(\mu\) cm och standardavvikelse
10 cm.
Säljaren av dessa plantor påstår att den genomsnittliga
höjden på plantorna är \(\mu=70\) cm.
Om \(\mu<70\) så har bonden rätt att underkänna
leveransen av plantor.

\begin{enumerate}[label=\alph*)]
\item Bonden misstänker att \(\mu<70\) cm.
För att testa detta, så mäter bonden 100 plantor
och finner att \(\bar{x}=68\) cm.
Kan bonden underkänna leveransen om han är villig
att ta en risk på 5\% att han gör fel?
Svara på denna fråga genom att ställa upp ett
lämpligt hypotestest (ställ upp hypotes \(H_0\)
och mothypotes \(H_1\)).
\hfill (3p)

\item Antag nu istället att bonden mäter 200 plantor
och finner att \(\bar{x}=69\).
Kan bonden underkänna leveransen om han är villig
att ta en risk på 5\% att han gör fel?
Svara på denna fråga genom att använda samma
hypotestest som du använde i a).
\hfill (3p)
\end{enumerate}

\textbf{English clarification:}
The notes mainly show the setup and calculation for
part a). Part b) says to use the same hypothesis test,
so we keep the same \(H_0\), \(H_1\), \(\alpha\),
and \(z\)-test formula, and only change \(n\)
and \(\bar{x}\).

\[
H_0:\mu=70,
\qquad
H_1:\mu<70.
\]

\[
\alpha=0.05,
\qquad
z_\alpha\approx -1.645.
\]

Now:

\[
n=200,
\qquad
\bar{x}=69.
\]

\begin{align*}
z
&=
\frac{69-70}
{10/\sqrt{200}} \\[2mm]
&=
\frac{-1}
{10/\sqrt{200}} \\[2mm]
&\approx
\frac{-1}{0.7071} \\[2mm]
&\approx -1.414.
\end{align*}

Compare with the critical value:

\[
-1.414>-1.645.
\]

So the test statistic is not in the rejection region.

Therefore:

\[
\text{Do not reject } H_0.
\]

The farmer cannot reject the delivery in part b).

\subsubsection*{Final answer}
\begin{enumerate}[label=\alph*)]
\item
\[
H_0:\mu=70,
\qquad
H_1:\mu<70.
\]

\[
z=-2.
\]

Since
\[
-2<-1.645,
\]
we reject \(H_0\).

The farmer can reject the delivery.

\item
Using the same test:

\[
z\approx -1.414.
\]

Since
\[
-1.414>-1.645,
\]
we do not reject \(H_0\).

The farmer cannot reject the delivery.
\end{enumerate}

\subsubsection*{What to remember}
For a left-tailed \(z\)-test:

\[
H_0:\mu=\mu_0,
\qquad
H_1:\mu<\mu_0.
\]

Use:

\[
z=
\frac{\bar{x}-\mu_0}
{\sigma/\sqrt{n}}.
\]

At 5\% left tail:

\[
z_\alpha\approx -1.645.
\]

Reject \(H_0\) if:

\[
z<-1.645.
\]

\end{document}
